\clearpage\pagestyle{otMain}\setlength{\otBlocEspacement}{2.5mm}\fontsize{11}{15}\selectfont\corrigesChapitre{17}{La symétrie, axes droits et obliques}
\begin{blocLivre}{cor-17-1}\begin{corrige}{1}{Combien d’axes ?}\demarche{Tester médianes et diagonales par pliage mental.}
Carré : $\res4$ ; rectangle : $\res2$ ; triangle équilatéral : $\res3$ ; cercle : $\res{\text{une infinité}}$ d’axes, tous ses diamètres.
\label{fig-cor-17-1}\figureGeo[corrige]{sym-1}
\end{corrige}\end{blocLivre}
\begin{blocLivre}{cor-17-2}\begin{corrige}{2}{Axe vertical}\demarche{Compter les carreaux jusqu’à l’axe.}
Chaque sommet traverse l’axe à même distance, sur sa ligne ; l’image, en vert, reproduit la figure retournée.
\label{fig-cor-17-2}\figureGeo[corrige]{sym-2}
\end{corrige}\end{blocLivre}
\begin{blocLivre}{cor-17-3}\begin{corrige}{3}{Axe horizontal}\demarche{Refléter colonne par colonne.}
Les points descendent sous l’axe à égale distance, colonne par colonne ; le pied du mât, posé sur l’axe, reste en place. L’image figure en vert.
\label{fig-cor-17-3}\figureGeo[corrige]{sym-3}
\end{corrige}\end{blocLivre}
\begin{blocLivre}{cor-17-4}\begin{corrige}{4}{Le point sur la diagonale}\demarche{Échanger « vers la droite » et « vers le haut ».}
Les déplacements s’échangent : le point à (3 ; 1) d’un point de l’axe va en (1 ; 3). Les trois images figurent en vert : $A'(1 ; 3)$, $B'(2 ; 5)$ et $C'(0 ; 4)$.
\label{fig-cor-17-4}\figureGeo[corrige]{sym-4}
\end{corrige}\end{blocLivre}
\begin{blocLivre}{cor-17-5}\begin{corrige}{5}{Segment et diagonale}\demarche{Traiter chaque extrémité puis joindre.}
A(2 ; 0) va en A'(0 ; 2) et B(3 ; 2) va en B'(2 ; 3), en comptant depuis le pied de l’axe : le segment image, vert, a la même longueur que [AB].
\label{fig-cor-17-5}\figureGeo[corrige]{sym-5}
\end{corrige}\end{blocLivre}
\begin{blocLivre}{cor-17-6}\begin{corrige}{6}{La figure oblique}\demarche{Associer les cases deux par deux à travers la diagonale.}
Chaque case colorée a sa jumelle obtenue en échangeant les deux directions le long de l’axe ; les cases ajoutées figurent en vert.
\label{fig-cor-17-6}\figureGeo[corrige]{sym-6}
\end{corrige}\end{blocLivre}
\begin{blocLivre}{cor-17-7}\begin{corrige}{7}{Chercher l’erreur}\demarche{Vérifier chaque sommet par l’échange des écarts.}
Le sommet $\res{C'}$ est placé en (2 ; 5) ; or C(5 ; 1) doit donner C'(1 ; 5), en échangeant ses deux coordonnées. Le point corrigé figure en vert.
\label{fig-cor-17-7}\figureGeo[corrige]{sym-7}
\end{corrige}\end{blocLivre}
\begin{blocLivre}{cor-17-8}\begin{corrige}{8}{Les lettres au miroir}\demarche{Tester le pliage, pas la rotation.}
Ces lettres se superposent à elles-mêmes après un demi-tour, pas par pliage : elles ont un centre de symétrie, pas un axe. Le pliage, lui, échoue : $\res{\text{aucun axe}}$. La symétrie axiale exige le pliage.
\label{fig-cor-17-8}\figureGeo[corrige]{sym-8}
\end{corrige}\end{blocLivre}
\begin{blocLivre}{cor-17-9}\begin{corrige}{9}{Deux axes perpendiculaires}\demarche{Enchaîner les deux reflets.}
On reflète d’abord à travers l’axe vertical, puis tout le résultat à travers l’horizontal : la figure finale montre 4 quarts identiques, ajoutés en vert.
\label{fig-cor-17-9}\figureGeo[corrige]{sym-9}
\end{corrige}\end{blocLivre}
\begin{blocLivre}{cor-17-10}\begin{corrige}{10}{Les longueurs conservées}\demarche{Citer la conservation des longueurs.}
La symétrie conserve les longueurs : l’image mesure $\res4$ carreaux, désormais comptés dans l’autre direction. La construction verte le confirme.
\label{fig-cor-17-10}\figureGeo[corrige]{sym-10}
\end{corrige}\end{blocLivre}
\begin{blocLivre}{cor-17-11}\begin{corrige}{11}{Le logo du club}\demarche{Compter les cases restantes après pliage mental.}
Chaque case colorée se reflète en échangeant ses écarts à l’axe : $\res5$ cases sont à colorier, en vert au corrigé.
\label{fig-cor-17-11}\figureGeo[corrige]{sym-11}
\end{corrige}\end{blocLivre}
\begin{blocLivre}{cor-17-12}\begin{corrige}{12}{Le pavage malin}\demarche{Appliquer l’axe le plus proche au dernier motif.}
Chaque axe renvoie le motif précédent retourné : le motif suivant, en vert, regarde dans l’autre sens. Les frises alternent ainsi endroit et envers.
\label{fig-cor-17-12}\figureGeo[corrige]{sym-12}
\end{corrige}\end{blocLivre}
\begin{blocLivre}{cor-17-13}\begin{corrige}{13}{Axe à trouver}\demarche{Chercher le pli qui superpose les moitiés.}
L’axe est la $\res{\text{diagonale}}$ qui échange les deux moitiés ; plié le long d’elle, le calque superpose exactement la figure. Tout autre pli déborde.
\label{fig-cor-17-13}\figureGeo[corrige]{sym-13}
\end{corrige}\end{blocLivre}
\begin{blocLivre}{cor-17-14}\begin{corrige}{14}{La rosace du vitrail}\demarche{Compter un quart puis multiplier par quatre.}
Chaque reflet multiplie les cases : le quart en compte 6, la rosace complète $6\times4=\res{24}$ cases, ajoutées en vert.
\label{fig-cor-17-14}\figureGeo[corrige]{sym-14}
\end{corrige}\end{blocLivre}
\begin{blocLivre}{cor-17-15}\begin{corrige}{15}{Vrai ou faux ?}\demarche{S’appuyer sur les propriétés de conservation.}
Première phrase : $\res{\text{faux}}$, la symétrie conserve longueurs et angles, l’image d’un carré est un carré superposable. Seconde : $\res{\text{vrai}}$, un point de l’axe ne bouge pas au pliage.
\label{fig-cor-17-15}\figureGeo[corrige]{sym-15}
\end{corrige}\end{blocLivre}
